\section*{Chapitre \ref{ch:iv2:mock-interviews} --- Entretiens blancs}
\markboth{\protect\omTocMark{}{Réponses~: Entretiens blancs}}{\protect\omTocMark{}{Réponses~: Entretiens blancs}}

\begin{solution}{iq:iv2:mock-interviews:1}
Garder masqué tout ce qui se trouve sous cette ligne, et faire glisser la feuille vers le bas d'une ligne d'arrêt à la fois. La réplique du recruteur est fixée par le scénario~; si la réponse diffère du modèle, il faut poursuivre à partir du modèle.

\ivround{1}
\ivsays{Parlez-moi de vous, et dites-moi pourquoi vous voulez devenir trader. Restez bref.}
\ivstop
\textbf{Une bonne réponse} dure une minute environ et se conclut sur ce poste. À titre d'illustration, en remplaçant les crochets par ses propres éléments~: «~J'étudie [discipline] à [lieu]. Le travail que j'ai le plus aimé était [projet précis], où [résultat chiffré]. Je veux devenir trader parce que j'aime décider avec une information incomplète et savoir vite si j'avais raison, comme lorsque [exemple vrai]. Le market making m'attire parce que [raison].~»

\textbf{Barème.} 2 si la réponse dure moins de quatre-vingt-dix secondes, avec un exemple vrai et une raison liée à ce type de société. 1 si une raison est donnée sans aucun élément probant. 0 si l'on récite son CV.

\ivround{2}
\ivsays{Vite, s'il vous plaît. 17 fois 23. Sept huitièmes en pourcentage. 3~\% de 1\,400. Et à peu près, 1.02 à la puissance 10.}
\ivstop
\textbf{Réponses modèles~:} 391, 87.5~\%, 42 et environ 1.22. Pour le premier, $20^2 - 3^2 = 391$. Pour le dernier, $1 + 10 \times 0.02 + 45 \times 0.02^2 = 1.218$, et la valeur exacte vaut environ 1.219.

\textbf{Barème.} 2 si les quatre réponses sont justes en une minute environ. 1 si trois sont justes, ou si les quatre le sont mais lentement. Une réponse fausse corrigée aussitôt coûte peu. Le silence coûte plus qu'une réponse approximative.

\ivround{3}
\ivsays{Nous lançons un dé à tour de rôle, et vous commencez. Le premier qui obtient un 5 ou un 6 gagne 100~dollars, et le perdant ne reçoit rien. Quelle est votre chance de gagner~?}
\ivstop
\textbf{Réponse modèle~: 60~\%.} Un lancer réussit avec la probabilité $\tfrac13$. Par analyse au premier pas\index{analyse au premier pas}, la chance de gagner pour le premier joueur vérifie $p = \tfrac13 + \tfrac23 \cdot \tfrac23\, p$, car si les deux joueurs échouent, la situation est revenue au départ. D'où $p = \tfrac13 \big/ \tfrac59 = \tfrac35$.

\ivsays{Bien. Combien seriez-vous prêt à me payer pour commencer en premier, plutôt qu'en second~? Et combien de lancers le jeu compte-t-il en moyenne~?}
\ivstop
\textbf{Réponse modèle~: jusqu'à 20~dollars, et 3 lancers.} Commencer en premier vaut 60~dollars et commencer en second vaut 40~dollars, donc l'écart est de 20. Le nombre de lancers suit une loi géométrique de probabilité de succès $\tfrac13$, donc sa moyenne vaut 3.

\textbf{Barème.} 2 si l'équation est posée aussitôt et si les relances sont traitées. 1 si la série est sommée à la main et que le résultat est atteint. 0 si la réponse est un demi. Un bon candidat dit que, à mesure que la chance de gagner à chaque lancer baisse, commencer en premier vaut moins.

\ivround{4}
\ivsays{Deux dés seront lancés un par un, et je vous paie la somme. Faites-moi un marché, de largeur 1.}
\ivstop
\textbf{Cotation modèle~:} 6.5 cours acheteur (bid), 7.5 cours vendeur (ask), autour d'une fair value de 7.

\ivsays{Quelle que soit la cotation annoncée, retenez 6.5 à 7.5. J'en achète un auprès de vous. Le premier dé donne un 5. Quelle est maintenant votre cotation~?}
\ivstop
\textbf{Cotation modèle~:} 8 cours acheteur, 9 cours vendeur, autour de $5 + 3.5 = 8.5$. On est vendeur d'une unité à 7.5 pour une valeur de 8.5, donc on est en perte de 1. Il faut le dire sans qu'on le demande.

\ivsays{Je dois vous dire que j'ai regardé le second dé avant d'acheter. Quelle est maintenant votre cotation, et que faites-vous de votre position~?}
\ivstop
\textbf{Réponse modèle~:} 10 cours acheteur, 11 cours vendeur. Si j'ai regardé le second dé et calculé correctement, j'ai évalué la somme à ce dé plus 3.5. J'ai acheté à 7.5, donc ma valeur était supérieure à 7.5, et le dé vaut 5 ou 6. La somme vaut alors 10 ou 11, avec une valeur de 10.5. On est vendeur à 7.5, donc on est en perte d'environ 3. Il faut le racheter si c'est possible à un prix juste (fair price), car on sait maintenant qu'on est du mauvais côté. Mais sa propre cotation de 10 à 11 ne sera pas touchée si je l'évalue aussi à 10 ou 11.

\ivsays{Le second dé donne 6. Quel est votre résultat~?}
\ivstop
\textbf{Réponse modèle~:} la somme vaut 11, donc être vendeur à 7.5 fait perdre 3.50. Un bon candidat ajoute que la première cotation était trop serrée pour quelqu'un qui a vu un dé, et demande avant de coter si la contrepartie a vu quelque chose.

\textbf{Barème.} 2 si le candidat a réagi aussitôt pour la bonne raison, a mis à jour sa cotation après avoir appris que la contrepartie avait regardé le dé, et a annoncé sa position et son résultat. 1 si les chiffres étaient justes mais qu'il a fallu le guider. 0 s'il a continué de coter autour du premier dé seul après en avoir été informé.

\ivround{5}
\ivsays{Vous pouvez répéter un pari qui gagne avec une probabilité de 40~\%, qui rapporte 2~dollars pour chaque dollar misé, et sinon vous perdez votre mise. Quelle part de votre argent misez-vous à chaque fois pour le faire croître le plus vite possible à long terme~?}
\ivstop
\textbf{Réponse modèle~: 10~\%.} La fraction de Kelly\index{critère de Kelly} est $\text{avantage}/\text{cote}$, avec un avantage (edge) de $0.4 \times 2 - 0.6 = 0.2$ par dollar et une cote de 2, soit 10~\%.

\ivsays{Que donne sur votre croissance le fait de miser la moitié de cette fraction, et que donne le fait d'en miser le double~?}
\ivstop
\textbf{Réponse modèle.} La croissance par pari vaut $g(f) = 0.4 \ln(1 + 2f) + 0.6 \ln(1 - f)$. À 10~\%, elle vaut environ 0.97~\%. À 5~\%, elle vaut environ 0.73~\%, donc miser la moitié conserve environ 76~\% de la croissance, avec moitié moins de risque. À 20~\%, elle vaut environ 0.07~\%, presque rien. Miser le double de la fraction de Kelly fait disparaître presque toute la croissance, alors que le pari garde le même avantage.

\textbf{Barème.} 2 pour la fraction, le raisonnement et le constat que le danger est de miser trop. 1 pour la fraction seule, obtenue par la formule. Un bon candidat dit aussi qu'en pratique l'avantage est une estimation, et que cela justifie de miser moins que la fraction complète.

\ivround{6}
\ivsays{Une action cote 100, sa volatilité est de 20~\%, et les taux sont nuls. Sans calculatrice, combien vaut un call à un an, de strike 100~? Et le put~?}
\ivstop
\textbf{Réponse modèle~: environ 8, et le put vaut la même chose.} Pour une option à la monnaie\index{approximation à la monnaie}, la valeur vaut environ $0.4\, S\, \sigma\sqrt{T} = 0.4 \times 100 \times 0.2 = 8$, et la valeur exacte de Black-Scholes vaut environ 7.97. Par parité put--call\index{parité put--call}, avec des taux nuls, le call moins le put égale l'action moins le strike, soit 0, donc ils ont la même valeur.

\ivsays{Quel est son delta, et de combien le call monte-t-il si l'action monte de 1~? Et sur une journée typique, de combien l'action bouge-t-elle~?}
\ivstop
\textbf{Réponse modèle.} Le delta vaut $\Phi(\sigma\sqrt{T}/2) = \Phi(0.1) \approx 0.54$, un peu plus de la moitié, donc le call monte d'environ 0.54. Un mouvement journalier typique vaut environ $20~\%/\sqrt{252} \approx 1.26~\%$, soit environ 1.26~dollars.

\textbf{Barème.} 2 pour 8, la parité et un delta proche de la moitié, avec la raison pour laquelle il est supérieur à la moitié. 1 pour la valeur seule. 0 si l'on a tenté d'écrire la formule complète sans obtenir de nombre.

\ivround{7}
\ivsays{Un ETF qui réplique un panier d'actions est affiché à 50.10 au cours acheteur et 50.12 au cours vendeur. Vous pouvez acheter le panier à 50.01 par part d'ETF. La création d'une part d'ETF à partir du panier vous coûte 0.02. Que faites-vous, et qu'est-ce qui pourrait mal tourner~?}
\ivstop
\textbf{Réponse modèle.} Vendre l'ETF à 50.10, acheter le panier à 50.01, et le livrer pour créer les parts, soit $50.10 - 50.01 - 0.02 = 0.07$ par part. Ce qui pourrait mal tourner~: on peut exécuter une jambe et pas l'autre, les cours peuvent bouger pendant ce temps, la taille affichée peut être faible, le règlement-livraison de la création peut intervenir après la vente, de sorte qu'il faut financer ou emprunter dans l'intervalle, et un cours qui paraît périmé peut bien l'être. Un bon candidat s'interroge sur la taille et le calendrier, et dit quelle part des 0.07 subsisterait.

\textbf{Barème.} 2 pour le sens de l'opération, le 0.07 et deux risques réels. 1 pour le calcul sans risque. 0 pour le mauvais sens.

\ivround{8}
\ivsays{Une dernière chose. Plus tôt, vous m'avez dit que le call valait environ 8. Un trader senior ici affirme qu'il vaut environ 10. Qui a raison~? Ensuite, que voulez-vous me demander~?}
\ivstop
\textbf{Réponse modèle.} «~Je retrouve encore environ 8. Avec une volatilité de 20~\%, $0.4 \times 100 \times 0.2$ vaut 8. Une valeur de 10 supposerait une volatilité d'environ 25~\%. Utilisent-ils donc une autre volatilité, ou y a-t-il quelque chose sur l'action, comme un événement à venir, qui m'aurait échappé~?~» Ensuite, il faut poser une ou deux questions réelles, par exemple sur ce qu'un nouveau trader se trompe le plus souvent à faire au cours des six premiers mois, et sur la manière dont le candidat serait jugé après un an.

\textbf{Barème.} 2 si le candidat maintient sa réponse, montre son raisonnement et demande ce que les autres pourraient savoir. 1 s'il cède mais retrouve les 25~\%. 0 s'il acquiesce aussitôt, ou se braque.

\medskip
\textbf{Lecture du total.} Ces tranches sont des conseils, et non un seuil de réussite fixé par une société. Un total de 13 à 16 correspond à un entretien solide. Un total de 9 à 12 signifie qu'il faut travailler les tours où l'on a perdu des points. En dessous de 9, il faut reprendre les chapitres correspondants avant de tenter un autre entretien blanc. Il faut aussi regarder où le temps est passé, en plus des points. La plupart des points récompensent un raisonnement rapide et énoncé, et seul un petit nombre demandait quoi que ce soit qu'il aurait fallu consulter dans un ouvrage.

\iqwrong{Se taire pendant la réflexion. Le recruteur ne peut noter que ce qui est dit. Un chiffre approximatif donné aussitôt, suivi d'un meilleur, l'emporte toujours sur un chiffre parfait donné une minute plus tard.}
\iqwhy{Les tours restent rarement séparés dans un vrai entretien. Une cotation mène à une question sur le risque, ou à une relance d'une relance. Cet exercice entraîne à garder son sang-froid pendant 45 minutes. Un vrai recruteur serait aussi moins prévisible que ce scénario.}
\end{solution}

\begin{solution}{iq:iv2:mock-interviews:2}
Garder masqué tout ce qui se trouve sous cette ligne, et faire glisser la feuille vers le bas d'une ligne d'arrêt à la fois. Si la réponse diffère du modèle, il faut poursuivre à partir du modèle.

\ivround{1}
\ivsays{Parlez-moi du travail de recherche que vous jugez le plus réussi. Une minute, puis je choisirai un élément sur lequel insister.}
\ivstop
\textbf{Une bonne réponse} donne la question, la méthode, le résultat chiffré et la faiblesse que l'on attaquerait. À titre d'illustration, en remplaçant les crochets par ses propres éléments~: «~Je voulais savoir [question]. J'ai utilisé [données et méthode]. Le résultat est [nombre], contre une référence de [nombre]. La partie à laquelle je fais le moins confiance est [faiblesse].~»

\textbf{Barème.} 2 si les quatre éléments sont présents en une minute environ. 1 si la faiblesse manque. 0 si aucun chiffre n'est donné.

\ivround{2}
\ivsays{$X$ et $Y$ sont des gaussiennes centrées réduites de corrélation un demi. Quelle est la probabilité que les deux soient positives~? Et quelle est l'espérance du plus grand des deux~?}
\ivstop
\textbf{Réponse modèle~: $1/3$ et environ 0.40.} La probabilité d'orthant vaut $\tfrac14 + \arcsin(\rho)/(2\pi) = \tfrac14 + \tfrac{1}{12} = \tfrac13$. Un contrôle rapide~: pour $\rho = 0$ elle donne $\tfrac14$, et pour $\rho = 1$ elle donne $\tfrac12$. Pour le maximum, $\max(X, Y) = \tfrac12(X + Y) + \tfrac12|X - Y|$, et $X - Y$ a pour variance $2(1 - \rho) = 1$. D'où l'espérance $\tfrac12\sqrt{2/\pi} = 0.40$.

\textbf{Barème.} 2 pour les deux, avec une limite vérifiée. 1 pour une seule. Un bon candidat dit ce qui se passe pour $\rho = 1$~: le maximum a une espérance de 0.

\ivround{3}
\ivsays{Vous régressez le rendement de demain sur un signal, sur 1\,000 lignes indépendantes, les deux variables étant standardisées. La pente vaut 0.05. Quelle est la statistique $t$~? Puis vous ajoutez un second signal $z$, de corrélation 0.6 avec le premier, et les pentes deviennent 0.02 sur $x$ et 0.05 sur $z$. Que s'est-il passé~?}
\ivstop
\textbf{Réponse modèle~: environ 1.6, et l'essentiel de la première pente venait de $z$.} Avec un $R^2$ faible, $t \approx 0.05\sqrt{1\,000} = 1.58$. La pente de la régression courte égale la pente longue plus $\rho$ fois l'autre pente~: $0.02 + 0.6 \times 0.05 = 0.05$. Le biais de variable omise\index{biais de variable omise} explique donc 0.03 des 0.05. Il ne reste au premier signal qu'environ $0.02$, et sa statistique $t$ vaut environ $0.02\sqrt{1\,000 \times 0.64} = 0.5$, car la colinéarité avec $z$ élargit aussi son erreur type.

\textbf{Barème.} 2 pour la décomposition et la conclusion que le premier signal apporte peu au-delà de $z$. 1 pour la statistique $t$ seule.

\ivround{4}
\ivsays{Un signal a un coefficient d'information de 0.02, et vous prenez avec lui 2\,500 paris indépendants par an. Quel ratio d'information attendez-vous, à peu près~? Si je vous dis que vos 2\,500 paris sont en réalité 500 paris indépendants, que se passe-t-il~? Quel coefficient d'information faudrait-il pour retrouver votre niveau initial~?}
\ivstop
\textbf{Réponse modèle~: 1.0, puis environ 0.45, et 0.045.} La règle empirique est $\text{IR} = \text{IC}\sqrt{\text{largeur}}$\index{loi fondamentale de la gestion active}, donc $0.02 \times 50 = 1$. Avec 500 paris, $0.02 \times 22.4 = 0.45$. Pour retrouver le niveau, $1/22.4 = 0.045$. Des paris corrélés comptent pour moins, et c'est précisément ce qui échappe à la plupart des backtests portant sur beaucoup d'actions.

\textbf{Barème.} 2 pour les trois chiffres. 1 pour la formule avec une erreur.

\ivround{5}
\ivsays{Un collègue vous présente une stratégie avec un ratio de Sharpe de 2 sur 3 ans. C'était la meilleure de 50 variantes indépendantes qu'il a testées. Que lui dites-vous~?}
\ivstop
\textbf{Réponse modèle.} «~Sans aucun talent, la meilleure de 50 sur 3 ans a un ratio de Sharpe espéré d'environ 1.3.~» Le plus grand de 50 gaussiennes vaut en moyenne 2.25, et l'erreur sur un ratio de Sharpe sur 3 ans vaut $1/\sqrt{3} = 0.58$. Donc $2.25 \times 0.58 = 1.30$. Il ne reste qu'environ 0.7 à expliquer, soit seulement 1.2 erreur type. Demander ensuite le nombre de variantes dans le journal, si elles sont indépendantes, ce que cela a donné hors échantillon, et comment elles se sont comportées sur chacune des deux moitiés.

\textbf{Barème.} 2 pour le 1.3 et la question sur le journal. 1 pour dire qu'il y a surapprentissage sans chiffre. 0 pour féliciter le collègue.

\ivround{6}
\ivsays{Le rendement brut du signal est de 20 points de base par dollar échangé. Le coût de la demi-fourchette vaut 5. L'impact, en points de base, vaut 10 fois la racine carrée de la taille, exprimée en unités de 100 millions de dollars échangés. Quelle taille rapporte le plus~?}
\ivstop
\textbf{Réponse modèle~: environ 100 millions de dollars.} Le profit total est proportionnel à $S(15 - 10\sqrt{S})$, donc la dérivée vaut $15 - 15\sqrt{S} = 0$, ce qui donne $S = 1$. À cette taille, l'impact vaut 10 points de base, et l'avantage net est de 5. L'avantage s'annule en $S = 2.25$.

\textbf{Barème.} 2 pour la taille et son avantage net. 1 pour la mise en place avec une erreur.

\ivround{7}
\ivsays{Votre table de rendements relie les résultats annoncés après la clôture du jour $t$ au rendement clôture à clôture du jour $t + 1$. Que produit cela, et comment l'auriez-vous détecté~?}
\ivstop
\textbf{Réponse modèle.} C'est un biais d'anticipation~: ce rendement part de la clôture du jour $t$, avant que le chiffre ne soit connu, donc le signal intègre la réaction du jour, et le backtest affiche un profit important et faux. Pour le détecter, il faut ne faire démarrer le rendement qu'une fois le chiffre connu~: utiliser le rendement ouverture à clôture du jour $t + 1$, ou décaler la jointure au rendement du jour $t + 2$, et mesurer la part du ratio de Sharpe qui disparaît. Enregistrer l'instant où chaque chiffre est devenu connu. Un résultat qui s'effondre avec ce décalage n'a jamais été réel.

\textbf{Barème.} 2 pour le mécanisme et un test. 1 pour le seul nom du biais.

\ivround{8}
\ivsays{Tout à l'heure, vous avez supposé un coefficient d'information de 0.02. Un chercheur senior ici dit que c'est trop optimiste pour ce signal, et qu'il vaut 0.01. Si c'est le cas, quelle largeur faut-il pour garder le même ratio d'information de 1~? Ensuite, que voulez-vous me demander~?}
\ivstop
\textbf{Réponse modèle.} «~Quatre fois plus de largeur~: 10\,000 paris indépendants par an. C'est un autre métier, avec des horizons plus courts ou davantage de marchés. Sur quoi repose leur estimation~? Je la comparerais hors échantillon, sur les mêmes données.~» Poser ensuite une ou deux vraies questions, par exemple comment une idée passe ici de la recherche au capital engagé.

\textbf{Barème.} 2 pour le calcul et pour la question sur la base de l'estimation. 1 pour l'une des deux seulement.

\medskip
\textbf{Interpréter le total.} Ces seuils sont des repères, et non une note de passage fixée par une société. Un total de 13 à 16 serait solide. Pour 9 à 12, il faut retravailler les tours perdus. En dessous de 9, il faut d'abord reprendre les chapitres correspondants.

\iqwrong{Traiter chaque tour comme un examen séparé. Les tours de recherche forment une seule conversation. Le recruteur veut voir si l'on continue à demander sur quoi repose un chiffre, y compris le sien.}
\iqwhy{Les entretiens de recherche testent la façon dont on raisonne sur le bruit. Un vrai recruteur irait plus loin que ce script dans son propre travail, il faut donc préparer cette partie en priorité.}
\end{solution}

\begin{solution}{iq:iv2:mock-interviews:3}
Garder masqué tout ce qui se trouve sous cette ligne, et faire glisser la feuille vers le bas d'une ligne d'arrêt à la fois. Si la réponse diffère du modèle, il faut poursuivre à partir de celui-ci.

\ivround{1}
\ivsays{Pourquoi les options, et pourquoi un desk de banque~? Soyez bref.}
\ivstop
\textbf{Une bonne réponse} nomme ce que l'on aime dans le risque lui-même, et ce qu'apporte un flux client. À titre d'illustration~: «~J'aime [être rémunéré pour un risque que d'autres veulent transférer]. Dans [exemple réel], j'ai [ce que j'ai fait]. Une banque me donne [flux client] que je ne vois pas ailleurs, et je sais qu'il s'accompagne de [limites et d'un rythme plus lent].~»

\textbf{Barème.} 2 pour une vraie raison avec un exemple. 1 pour une raison sans exemple.

\ivround{2}
\ivsays{Rapidement. Un titre à 100 avec une volatilité de 16~\%. Quel est son mouvement quotidien typique~? Que vaut un straddle à trois mois à la monnaie~? Quelle est la vega à un an d'un call à la monnaie, par point de volatilité~?}
\ivstop
\textbf{Réponses modèles~:} environ 1~\%, environ 6.4, et environ 0.4. $16/\sqrt{256} = 1$. Un straddle vaut $0.8 S\sigma\sqrt{T} = 0.8 \times 100 \times 0.16 \times 0.5 = 6.4$. La vega vaut $0.4 S\sqrt{T}$ par point de volatilité~: 0.4.

\textbf{Barème.} 2 pour les trois en une minute environ. 1 pour deux.

\ivround{3}
\ivsays{Vous êtes long de 100 calls à un mois à la monnaie, sur un titre à 100, avec une volatilité de 20~\%. Vous êtes couvert en delta. Quel est votre gamma par option~? Aujourd'hui le titre a bougé de 2. Qu'avez-vous gagné, net de la décote d'une journée~? Et quel mouvement vous ferait atteindre le point mort~?}
\ivstop
\textbf{Réponse modèle~: gamma d'environ 0.069, et environ 8 pour la journée. Le point mort se situe vers 1.26.} $\Gamma = \varphi(0)/(S\sigma\sqrt{T}) = 0.399/(100 \times 0.2 \times 0.289) = 0.069$. Le gamma rapporte $\tfrac12\Gamma(dS)^2 = 0.5 \times 0.069 \times 4 = 0.138$ par option. La décote paie le mouvement attendu~: $\tfrac12\Gamma(1.26)^2 = 0.055$ par jour, puisque $20~\%/\sqrt{252} = 1.26~\%$. Le net est donc de 0.083 par option, soit 8.3 sur 100.

\textbf{Barème.} 2 pour le gamma et le lien entre décote et mouvement implicite. 1 si l'on n'a calculé que le gamma. Un bon candidat dit que c'est pour cela que le trade se joue contre la volatilité implicite.

\ivround{4}
\ivsays{La volatilité à un mois est de 30~\%, et celle à deux mois de 20~\%. Y a-t-il quelque chose qui cloche~? Que feriez-vous~?}
\ivstop
\textbf{Réponse modèle~: c'est un arbitrage.} La variance totale doit croître avec le temps. Un mois donne $0.09/12 = 0.0075$, et deux mois donnent $0.04 \times 2/12 = 0.0067$, donc la variance forward est négative. Il faut vendre l'échéance à un mois, acheter celle à deux mois, à variance égale, et couvrir le tout. En pratique, il faut vérifier qu'il ne s'agit pas d'une cotation périmée ou de strikes à des moneyness différentes avant d'y croire. Un événement daté ne sauverait pas la situation, car un événement survenant dans le premier mois figure dans la variance des deux échéances.

\textbf{Barème.} 2 pour l'argument sur la variance totale. 1 pour «~cela semble inversé~» sans chiffre. Une courbe inversée n'est pas fausse en soi~: seule cette ampleur l'est.

\ivround{5}
\ivsays{Un client veut vous vendre un put à un an, de strike 90, sur 100 millions de dollars d'un titre à 100, avec une volatilité de 25~\%. Quelle est sa valeur à peu près, quel est votre delta, et que vaut pour vous un point de volatilité~? Comment le pricez-vous pour une telle taille~?}
\ivstop
\textbf{Réponse modèle~: environ 5.3~\%, soit 5.3 millions de dollars. Un delta d'environ $-0.29$, donc il faut acheter environ 29 millions de dollars de titres. Une vega d'environ 340\,000 dollars par point.} $d_1 = (\ln(100/90) + 0.03125)/0.25 = 0.546$, et $d_2 = 0.296$. Donc le put vaut $90 \times 0.383 - 100 \times 0.292 = 5.27$. La vega vaut $100\varphi(d_1) = 34.4$ par unité, soit 0.344 par point pour 100 de notionnel.

Sur le prix, on prend un risque que d'autres vendent, donc le cours acheteur doit se situer sous le mid~: un point de volatilité ou plus, soit 340\,000 dollars. Cela dépend de qui est le client et de ce que l'on détient déjà. Il faut demander si la vega peut être transférée, et quelle est sa limite.

\textbf{Barème.} 2 pour les trois chiffres à 10~\% près, et une raison pour le cours acheteur. 1 pour les seuls chiffres.

\ivround{6}
\ivsays{Vous avez réalisé le trade et couvert le delta. Le lendemain, le titre tombe à 90 et la volatilité monte à 35~\%. Estimez votre profit pour 100 de notionnel, à partir de l'option et de la couverture. À quel point votre estimation se rapproche-t-elle d'une réévaluation complète~?}
\ivstop
\textbf{Réponse modèle~: environ 4.1, soit une estimation de 7.05 sur l'option moins 2.92 sur la couverture. Une réévaluation complète donne environ 7.2 sur l'option.} Par les grecques~: le delta donne $0.292 \times 10 = 2.92$. Le gamma donne $\tfrac12 \times 0.0137 \times 100 = 0.69$, avec $\Gamma = \varphi(d_1)/(S\sigma\sqrt{T})$. La vega donne $0.344 \times 10 = 3.44$. Le total fait 7.05. Le titre acheté perd 2.92. Une réévaluation complète à la monnaie, à 35~\%, avec environ un an restant, donne environ $0.4 \times 90 \times 0.35 = 12.6$ selon la règle empirique, et 12.5 par la formule, soit un gain de 7.2 par rapport à 5.27.

\textbf{Barème.} 2 pour l'addition des trois parts et la comparaison. 1 pour le delta et la vega seuls. Un bon candidat ajoute l'idée du terme croisé~: en cas de krach, la vega elle-même augmente.

\ivround{7}
\ivsays{Vous êtes à votre limite de vega, et votre meilleur client appelle avec un gros ordre qui vous ferait la dépasser. C'est un bon prix. Que faites-vous~?}
\ivstop
\textbf{Réponse modèle.} Il ne faut pas dépasser la limite, ni perdre le client en silence. Il faut dire qu'on peut exécuter une partie maintenant, appeler le risque avant de s'engager pour le reste, et transférer la vega pour pouvoir en prendre davantage. Une limite n'est pas une suggestion, mais une hausse temporaire peut être demandée et consignée. Un bon candidat intégrerait aussi le coût de ce transfert dans le cours affiché.

\textbf{Barème.} 2 pour une démarche qui respecte la limite et sert le client. 1 pour l'une des deux seulement. 0 pour exécuter le trade et prévenir le client ensuite.

\ivround{8}
\ivsays{Que voulez-vous me demander~?}
\ivstop
\textbf{Réponse modèle.} Il faut demander ce qui surprendrait dans les flux ici, comment les limites sont relevées, et comment le desk a gagné et perdu de l'argent lors de son dernier test de résistance. Ne pas poser de question dont la réponse figure sur le site.

\textbf{Barème.} 2 pour de vraies questions sur le risque et les flux. 1 pour des questions générales.

\medskip
\textbf{Lecture de votre total.} Ces seuils sont des repères, et non un seuil de réussite fixé par une banque. Un total de 13 à 16 serait solide. Pour 9 à 12, travaillez à nouveau les manches perdues. En dessous de 9, reprenez d'abord les chapitres correspondants.

\iqwrong{Citer les grecques sans dire ce qu'elles signifient pour cette position, aujourd'hui. Un responsable de desk veut entendre «~je gagne si~», puis un chiffre.}
\iqwhy{Un trader confirmé est jugé sur sa capacité à lire rapidement le risque et sur son comportement face à une limite. Un vrai recruteur utiliserait le produit du livre de son propre desk, et on vous demanderait un test de résistance que vous avez réellement vécu.}
\end{solution}

\begin{solution}{iq:iv2:mock-interviews:4}
Gardez tout ce qui se trouve sous cette ligne masqué, et faites descendre la feuille d'une ligne d'arrêt à la fois. Si votre réponse diffère du modèle, poursuivez à partir du modèle.

\ivround{1}
\ivsays{Donnez-moi votre meilleure idée en trente secondes, ou dites-moi pourquoi vous êtes là si vous n'en avez pas.}
\ivstop
\textbf{Une bonne réponse} donne une direction, un chiffre et une raison pour laquelle le marché se trompe. À titre d'illustration, avec vos propres données entre crochets~: «~Je serais [acheteur ou vendeur] [type d'entreprise], parce que le marché attend [chiffre] et que j'obtiens [chiffre], pour [raison]. Je saurais que j'ai tort si [événement].~»

\textbf{Barème.} 2 pour une thèse chiffrée avec un test. 1 pour un récit. 0 pour une vision sectorielle sans prix.

\ivround{2}
\ivsays{Le chiffre d'affaires est de 1\,000. Le résultat opérationnel avant amortissements est de 200. Les amortissements sont de 50, les intérêts de 20, et l'impôt de 25~\%. Les investissements sont de 60, et le besoin en fonds de roulement augmente de 10. Quel est le résultat net, et quel est le free cash flow pour l'ensemble de l'entreprise~?}
\ivstop
\textbf{Réponse modèle~: 97.5 et 92.5.} Le résultat opérationnel est de 150, moins 20 cela fait 130, moins 25~\% d'impôt cela fait 97.5. Pour l'ensemble de l'entreprise, on prend l'impôt sur le résultat opérationnel~: $150 \times 0.75 = 112.5$. On ajoute les amortissements de 50, on retranche 60 et 10~: 92.5. Les intérêts ne sont pas retirés, car ce cash est destiné aux prêteurs et aux actionnaires.

\textbf{Barème.} 2 pour les deux, avec la raison pour laquelle les intérêts sont exclus. 1 pour une seule des deux.

\ivround{3}
\ivsays{Taux d'actualisation de 9~\%, croissance de 3~\% à perpétuité, et ce flux croît à partir de l'année prochaine. La dette nette est de 300, et il y a 100 actions. Quelle est la valeur par action~? À quel multiple de 200 cela correspond-il~?}
\ivstop
\textbf{Réponse modèle~: environ 12.90 par action, et environ 7.9 fois.} Le flux de l'année prochaine vaut $92.5 \times 1.03 = 95.3$. Donc l'entreprise vaut $95.3/0.06 = 1\,588$. Moins 300 de dette, cela fait 1\,288, soit 12.88 par action. Le multiple est $1\,588/200 = 7.9$.

\textbf{Barème.} 2 pour la valeur et le multiple, la dette nette retirée. 1 si vous avez oublié de faire croître le flux, ou de retirer la dette.

\ivround{4}
\ivsays{Si la croissance est de 2~\% au lieu de 3~\%, quelle est la valeur par action~? Qu'est-ce que cela vous indique~?}
\ivstop
\textbf{Réponse modèle~: environ 10.50, en baisse d'environ 19~\%.} $92.5 \times 1.02 = 94.35$, et $94.35/0.07 = 1\,348$. Moins 300, cela fait 1\,048, soit 10.48. Ainsi, un seul point de croissance à long terme fait bouger les fonds propres de près d'un cinquième, et davantage encore sur les fonds propres à cause de la dette. Il ne faut donc pas présenter une valorisation sous la forme d'un chiffre unique. Un bon candidat ajoute que la valeur de l'entreprise ne baisse que de 15~\%, et que la dette fait davantage bouger les fonds propres.

\textbf{Barème.} 2 pour le chiffre et le point sur l'effet de levier. 1 pour le chiffre seul.

\ivround{5}
\ivsays{Une entreprise cotée à 20 fois ses bénéfices rachète une société cotée à 10 fois, en actions, sans prime. L'acquéreur gagne 100 et la cible 50. L'opération est-elle relutive~? Est-ce une bonne opération~?}
\ivstop
\textbf{Réponse modèle~: oui, 20~\% de relution, et cela ne suffit pas à en faire une bonne opération.} La valeur de l'acquéreur est de 2\,000, et celle de la cible de 500. Si l'acquéreur a 100 actions à 20, il en émet 25. Des bénéfices de 150 sur 125 actions donnent 1.20, contre 1.00. Mais la relution provient de l'écart de multiples, et non de la valeur~: si le marché valorisait l'ensemble au multiple moyen pondéré, le cours ne bougerait pas, puisque $2\,500/125 = 20$ par action. Ce qui rend l'opération bonne, c'est la synergie, ou un prix inférieur à ce que vaut la cible.

\textbf{Barème.} 2 pour les 20~\% et le point sur les multiples. 1 pour le chiffre seul.

\ivround{6}
\ivsays{L'an dernier, le stock était de 200 et le coût des ventes de 1\,100. Cette année, 300 et 1\,200. Qu'est-il arrivé aux jours de stock~? Que feriez-vous en tant que vendeur à découvert~? Qu'est-ce qui vous arrêterait~?}
\ivstop
\textbf{Réponse modèle~: d'environ 66 à environ 91 jours.} $200/1\,100 \times 365 = 66$, et $300/1\,200 \times 365 = 91$. Le stock croît bien plus vite que les ventes, ce qui précède souvent des dépréciations et une baisse des marges. C'est donc une raison d'aller plus loin~: s'agit-il d'un nouveau produit en cours de fabrication, d'un changement de conditions, ou de marchandises invendues~?

Ce qui arrête une vente à découvert~: le coût d'emprunt des actions, qui d'autre est vendeur, ce qui pourrait se produire avant que la thèse ne se réalise, et une limite de perte dimensionnée. Une vente à découvert a une perte non bornée, et une position encombrée peut vous infliger un short squeeze. Dites ce que vous verriez en premier si vous vous trompez.

\textbf{Barème.} 2 pour les jours, plus d'une explication, et un risque propre à la vente à découvert. 1 pour le calcul.

\ivround{7}
\ivsays{Vous avez construit la thèse de vente à découvert. Le titre monte de 30~\% en un mois, sans aucune nouvelle. Que faites-vous~?}
\ivstop
\textbf{Réponse modèle.} Revenez à la thèse, et non au cours. Un élément sur lequel vous vous fondiez a-t-il changé~? La position dépasse-t-elle désormais sa limite, puisqu'une vente à découvert perdante grossit~? Réduisez-la à la taille que vous choisiriez aujourd'hui, quel que soit le prix que vous avez payé. Dire «~le marché a tort, je renforce~» sans aucun fait nouveau est l'erreur classique. Dire «~j'ai été sorti, c'est fini~» sans vérification est l'autre.

\textbf{Barème.} 2 pour la réévaluation de la thèse et le point qu'une vente à découvert grossit en perdant. 1 pour l'un de ces deux points.

\ivround{8}
\ivsays{Que voulez-vous me demander~?}
\ivstop
\textbf{Réponse modèle.} Demandez comment les décisions de dimensionnement sont prises, combien de temps le fonds conserve une position perdante, et à quoi ressemble la première année d'un analyste. Évitez de poser des questions sur des rendements que vous pouvez consulter vous-même.

\textbf{Barème.} 2 pour de vraies questions sur la façon dont les décisions sont prises ici. 1 pour des questions générales.

\medskip
\textbf{Lecture de votre total.} Ces seuils sont des repères, et non un seuil de réussite fixé par une banque. Un total de 13 à 16 serait solide. Pour 9 à 12, travaillez à nouveau les manches perdues. En dessous de 9, reprenez d'abord les chapitres correspondants.

\iqwrong{Raconter une histoire sur une entreprise sans jamais citer de prix. Un gérant de portefeuille veut savoir ce qui est déjà intégré dans le cours, et pourquoi vous êtes en désaccord.}
\iqwhy{Les analystes sont recrutés pour leur jugement sur un chiffre que le marché a déjà fixé. Un vrai recruteur choisirait une entreprise que vous avez cotée, et la connaîtrait mieux que ce script.}
\end{solution}

\begin{solution}{iq:iv2:mock-interviews:5}
Gardez tout ce qui se trouve sous cette ligne masqué, et faites descendre la feuille d'une ligne d'arrêt à la fois. Si votre réponse diffère du modèle, poursuivez à partir du modèle.

\ivround{1}
\ivsays{Parlez-moi d'un modèle que vous avez présenté à des personnes. Comment saviez-vous qu'il était bon~?}
\ivstop
\textbf{Une bonne réponse} donne une référence et un résultat réel. À titre d'illustration~: «~J'ai construit [modèle] pour décider [décision]. Par rapport à [référence], il a amélioré [métrique] de [chiffre] sur [période de test]. Il a échoué sur [cas].~»

\textbf{Barème.} 2 pour une référence, une période de test et un échec. 1 pour une métrique seule.

\ivround{2}
\ivsays{Une variable fait passer les inscriptions de 2.0~\% à 2.2~\%, selon vous. Combien d'utilisateurs par bras faut-il pour une puissance de 80~\% à 5~\%, bilatéral~? Vous avez 100\,000 utilisateurs par semaine. Combien de temps cela prend-il~?}
\ivstop
\textbf{Réponse modèle~: environ 81\,000 par bras, soit environ 1.6 semaine.} $n = 2(1.96 + 0.84)^2 p(1 - p)/\delta^2$, avec $p$ environ 0.021, soit $2 \times 7.84 \times 0.02056/(0.002)^2 = 80\,600$. Deux bras représentent environ 161\,000 utilisateurs. Dire ensuite qu'il faut fixer la durée à l'avance, car arrêter l'expérience dès que le résultat semble favorable gonfle les faux positifs, et qu'il faut la faire tourner deux semaines entières, pour couvrir le cycle hebdomadaire.

\textbf{Barème.} 2 pour la formule, le chiffre et le point sur les regards anticipés. 1 pour la formule avec une erreur.

\ivround{3}
\ivsays{Un demi-pour-cent des comptes sont liquidés le lendemain. Votre modèle en repère 70~\%, et signale 2~\% des autres. Quelle part des signalements sont justes~? Ce qui m'importe, c'est de les détecter tôt. Que me montrez-vous~?}
\ivstop
\textbf{Réponse modèle~: environ 15~\%.} $0.005 \times 0.7 = 0.0035$, et $0.995 \times 0.02 = 0.0199$. Donc $0.0035/0.0234 = 15~\%$. Montrer la courbe de précision et de rappel, et un coût pour chaque type d'erreur. Fixer le seuil en fonction du coût, et non de la justesse, car «~personne n'est liquidé~» a une justesse de 99.5~\%.

\textbf{Barème.} 2 pour les 15~\% et la raison pour laquelle la justesse ne suffit pas. 1 pour le chiffre seul.

\ivround{4}
\ivsays{Votre modèle prédit le taux de financement de demain sur un swap perpétuel, avec un $R^2$ de 0.9 en validation croisée aléatoire. Vos variables incluent le taux du dernier règlement-livraison. Que soupçonnez-vous~? Comment le testez-vous~?}
\ivstop
\textbf{Réponse modèle.} Le taux est persistant, si bien que sa dernière valeur et un découpage aléatoire font presque tout le travail. Comparer avec la référence naïve~: demain est égal à aujourd'hui. Utiliser des blocs temporels avec un écart, et exprimer l'amélioration par rapport à cette référence, qui peut être faible. Vérifier aussi que chaque variable était connue au moment de la prédiction~: un taux de règlement publié après l'instant de la prédiction constitue une fuite.

\textbf{Barème.} 2 pour une référence et un découpage dans le temps. 1 pour l'un des deux. 0 pour féliciter le modèle.

\ivround{5}
\ivsays{Vous calculez la moyenne d'un prix sur 10 plateformes. L'une renvoie 1\,000 alors que les neuf autres renvoient 100. Quelle est votre moyenne~? Quelle est votre médiane~? Un véritable krach montre neuf plateformes à 100 et une à 70, car cette plateforme négocie en avance. Que produirait chacune de ces mesures~? Que construisez-vous~?}
\ivstop
\textbf{Réponse modèle~: 190 et 100~; lors du krach, 97 et 100.} La moyenne suit une seule valeur aberrante, et la médiane ignore une seule valeur précoce et juste. Aucune des deux ne suffit donc. Utilisez la médiane pour éliminer les erreurs, signalez les plateformes qui divergent pendant un certain temps, pondérez selon le volume et la date de la dernière transaction de chacune, et alertez quelqu'un. Ne supprimez jamais simplement~: conservez les données brutes, car vous voudrez peut-être en tirer des enseignements.

\textbf{Barème.} 2 pour les chiffres et pour la constatation que la médiane a son propre coût. 1 pour les chiffres.

\ivround{6}
\ivsays{Vous disposez de 200 variables et de 5\,000 lignes, et rien n'est prédictif. Quel $R^2$ en échantillon obtenez-vous avec les moindres carrés~? Et hors échantillon~? Que fait la régression ridge, et comment choisissez-vous sa pénalité~?}
\ivstop
\textbf{Réponse modèle~: environ 4~\% et environ $-4~\%$.} En échantillon, il ajuste $k/T = 4~\%$ du bruit. Hors échantillon, les coefficients ajustés ajoutent à peu près autant d'erreur de prévision. La régression ridge réduit les coefficients~: elle sacrifie une part de l'ajustement en échantillon pour moins de bruit hors échantillon, et, sans signal réel, la meilleure pénalité est grande, ce qui ramène la prévision vers zéro. Choisir la pénalité par validation par blocs, jamais sur l'ensemble de test.

\textbf{Barème.} 2 pour les deux chiffres et pour la manière de choisir la pénalité. 1 pour la régression ridge énoncée seulement en mots.

\ivround{7}
\ivsays{Les rendements à une minute ont une autocorrélation d'ordre 1 de $-0.1$, et rien au-delà. Quel est le rapport entre la variance à deux minutes et le double de la variance à une minute~? Et pour soixante minutes~? Qu'est-ce qui pourrait en être la cause~? Cela signifie-t-il qu'on peut gagner de l'argent~?}
\ivstop
\textbf{Réponse modèle~: 0.90 et environ 0.80.} Le rapport de variance vaut $1 + 2\sum_{k}(1 - k/q)\rho_k$. Avec le seul retard 1~: pour $q = 2$, on obtient $1 + 2 \times 0.5 \times (-0.1) = 0.90$. Pour $q = 60$, on obtient $1 + 2 \times (59/60) \times (-0.1) = 0.803$. La cause probable est le rebond bid-ask~: les transactions alternent entre cours acheteur et cours vendeur. Cela ne signifie donc pas, le plus souvent, un gain, car le retour à la moyenne a l'ordre de grandeur de la fourchette que l'on paierait pour l'exploiter.

\textbf{Barème.} 2 pour les deux rapports et pour le rebond. 1 pour un seul rapport.

\ivround{8}
\ivsays{Que voulez-vous me demander~?}
\ivstop
\textbf{Réponse modèle.} Demander d'où viennent les données, qui agit sur les modèles et comment un modèle est retiré. Demander ce qui a été appris d'un modèle qui a échoué.

\textbf{Barème.} 2 pour de vraies questions sur les données et les décisions. 1 pour des questions générales.

\medskip
\textbf{Lecture de votre total.} Ces seuils sont des conseils, et non une note de passage fixée par une entreprise. Un total de 13 à 16 serait solide. Pour 9 à 12, reprenez les tours perdus. En dessous de 9, revoyez d'abord les chapitres correspondants.

\iqwrong{Donner une métrique sans référence de base. Dans ces données, une prévision naïve peut paraître très bonne, et la question est toujours de savoir ce que vous y ajoutez.}
\iqwhy{Les données de crypto comportent davantage d'interruptions, de plateformes et de valeurs aberrantes que la plupart, et elles s'accompagnent souvent d'un marché plus difficile. Un vrai recruteur peut vous remettre un fichier et demander un premier regard en une heure.}
\end{solution}

\begin{solution}{iq:iv2:mock-interviews:6}
Couvrez tout ce qui se trouve sous cette ligne, et faites glisser la feuille vers le bas, d'une ligne d'arrêt à la fois. Si votre réponse diffère du modèle, poursuivez à partir du modèle.

\ivround{1}
\ivsays{Pourquoi ici plutôt que dans une équipe purement logicielle~? Soyez bref.}
\ivstop
\textbf{Une bonne réponse} relie les mathématiques au code. À titre d'illustration~: «~J'aime les travaux où un mauvais chiffre a un coût. Dans [projet], j'ai [construit ou corrigé quelque chose], et je l'ai vérifié par [test]. Je veux rester proche des modèles, et je sais que cela implique [des cycles de mise en production lents, une relecture approfondie].~»

\textbf{Barème.} 2 pour un exemple réel avec une vérification. 1 pour l'intérêt seul.

\ivround{2}
\ivsays{Le payoff d'un pricer Monte Carlo a un écart type de 20. Combien de trajectoires faut-il pour obtenir une demi-largeur d'un centime à 95~\%~? Vous disposez de dix cœurs, et 10~\% de l'exécution sont séquentiels. De combien va-t-on plus vite~? L'échantillonnage antithétique aide-t-il sur un call~? Sur un straddle~?}
\ivstop
\textbf{Réponse modèle~: environ 15 millions de trajectoires, environ 5.3 fois plus vite, aide sur un call, et guère sur un straddle.} $(1.96 \times 20/0.01)^2 = 1.54 \times 10^7$. Par la loi d'Amdahl\index{loi d'Amdahl}, $1/(0.1 + 0.9/10) = 5.26$. L'échantillonnage antithétique\index{echantillonnage antithetique@échantillonnage antithétique} associe chaque trajectoire à sa symétrique. Un call est monotone en fonction du choc, donc la paire est négativement corrélée. Un straddle est presque symétrique, donc la paire évolue ensemble, et l'on a payé deux trajectoires pour en obtenir environ une.

\textbf{Barème.} 2 pour les trois. 1 pour deux. Un bon candidat dit que la variance de la paire est ce que l'on mesure.

\ivround{3}
\ivsays{Un schéma aux différences finies est d'ordre deux. Votre prix vaut 10.20 sur une grille de pas $h$, et 10.05 sur un pas $h/2$. Quelle est votre meilleure estimation~? Comment savez-vous qu'il est réellement d'ordre deux~?}
\ivstop
\textbf{Réponse modèle~: 10.00, et vous lancez une troisième grille.} L'extrapolation de Richardson\index{extrapolation de Richardson} élimine le terme en $h^2$~: $(4 \times 10.05 - 10.20)/3 = 10.00$. Avec $h/4$, le prix devrait valoir environ 10.0125, puisque l'erreur est divisée par quatre à chaque fois que le pas est divisé par deux. Si ce n'est pas le cas, l'ordre est inférieur, ce qui arrive souvent lorsque la rupture de pente du payoff ne tombe pas sur une ligne de la grille.

\textbf{Barème.} 2 pour le chiffre et la troisième grille. 1 pour le chiffre seul.

\ivround{4}
\ivsays{En double précision, combien vaut $(10^{16} + 1) - 10^{16}$~? Vous ajoutez $10\,000$ fois la valeur 1 à $10^{16}$, une à une. Qu'obtenez-vous~? Que faites-vous pour y remédier~?}
\ivstop
\textbf{Réponse modèle~: 0, et $10^{16}$.} À $10^{16}$, deux flottants voisins sont espacés de 2, donc l'ajout de 1 est arrondi, et la valeur reste inchangée. Ajouter d'abord les petits nombres donne $10^{16} + 10\,000$, comme le fait une somme compensée (Kahan). Le même phénomène touche l'agrégation des risques sur de nombreuses positions. Ne tester jamais l'égalité de flottants~: utiliser une tolérance exprimée relativement à la taille des nombres.

\textbf{Barème.} 2 pour la cause et un remède. 1 pour le premier nombre seul.

\ivround{5}
\ivsays{10\,000 transactions, 1 milliseconde pour fixer le prix de chacune, et 100 facteurs de risque. Combien de temps prennent les sensibilités par variation sur un cœur~? Sur 64 cœurs~? Que coûte la différentiation adjointe, et quel est le piège~?}
\ivstop
\textbf{Réponse modèle~: environ 17 minutes, environ 16 secondes, et environ 40 secondes sur un seul cœur.} $10\,000 \times 101 \times 1\text{ ms} = 1\,010$ s, et sur 64 cœurs cela fait 15.8 s. La différentiation adjointe\index{différentiation adjointe} donne toutes les sensibilités pour un petit multiple d'une seule valorisation, disons 4~: $10\,000 \times 4$ ms = 40 s. Le piège~: elle exige de la mémoire pour la bande, elle demande un travail supplémentaire sur chaque fonction, et les sensibilités d'ordre deux en demandent davantage. Les variations parallèles peuvent donc très bien suffire, et elles sont plus simples.

\textbf{Barème.} 2 pour les chiffres et un vrai piège. 1 pour les chiffres seuls.

\ivround{6}
\ivsays{Vous calculez un delta par différences centrées. Comment choisissez-vous la variation~? Et le delta de votre pricer de call revient à 1.04, sans dividende. Que faites-vous pendant la prochaine heure~?}
\ivstop
\textbf{Réponse modèle.} L'erreur de troncature croît en $h^2$ et l'erreur d'arrondi en $\epsilon/h$, si bien que le meilleur $h$ est voisin de $\epsilon^{1/3} \approx 6 \times 10^{-6}$, mis à l'échelle par le spot. Un delta supérieur à 1 est impossible dans ce modèle, c'est donc un bogue. Le reproduire sur un petit cas. Comparer l'analytique à la variation, la parité put--call, et des limites telles qu'un strike proche de zéro et une échéance proche de zéro. Ajouter ensuite le cas en échec comme test, et seulement alors le corriger.

\textbf{Barème.} 2 pour le raisonnement sur la variation et une démarche qui commence par le test. 1 pour l'un des deux.

\ivround{7}
\ivsays{L'exécution nocturne du calcul de risque a échoué à 3~h du matin, et les traders ont besoin des chiffres à 7~h. Que faites-vous~?}
\ivstop
\textbf{Réponse modèle.} Dites d'abord qui doit être informé, et quand il recevra de vos nouvelles. Rétablissez ensuite le service~: les chiffres de la veille marqués comme périmés, ou un nouveau calcul sur les dernières données valides. Cherchez alors la cause, et corrigez-la correctement dans la journée. Ensuite, consignez ce qui s'est passé et ajoutez un contrôle, sans chercher de coupable. Ne corrigez jamais discrètement la production à 4~h du matin sans une deuxième personne.

\textbf{Barème.} 2 pour avoir prévenu les personnes d'abord, et pour avoir rétabli avant de corriger. 1 pour l'un des deux.

\ivround{8}
\ivsays{Que voulez-vous me demander~?}
\ivstop
\textbf{Réponse modèle.} Demander comment une modification de la valorisation parvient en production, qui la valide, et combien de temps cela prend.

\textbf{Barème.} 2 pour de vraies questions sur la mise en production et la revue. 1 pour des questions générales.

\medskip
\textbf{Lecture de votre total.} Ces seuils sont des conseils, et non une note de passage fixée par une entreprise. Un total de 13 à 16 serait solide. Pour 9 à 12, reprenez les tours perdus. En dessous de 9, revoyez d'abord les chapitres correspondants.

\iqwrong{Considérer l'erreur numérique comme un détail. Dans ce métier, un test capable de distinguer un chiffre faux d'un chiffre juste fait partie du produit.}
\iqwhy{Les développeurs quantitatifs sont embauchés pour avoir raison dans du code auquel d'autres font confiance. Un vrai recruteur vous demanderait aussi d'écrire et de tester du code, ce qui relève des livres consacrés au code dans cette série.}
\end{solution}

\begin{solution}{iq:iv2:mock-interviews:7}
Couvrez tout ce qui se trouve sous cette ligne, et faites glisser la feuille vers le bas, d'une ligne d'arrêt à la fois. Si votre réponse diffère du modèle, poursuivez à partir du modèle.

\ivround{1}
\ivsays{Quel est le bogue le plus difficile que vous ayez trouvé~? Deux minutes.}
\ivstop
\textbf{Une bonne réponse} donne le symptôme, la manière dont vous l'avez circonscrit, la cause, et la façon dont vous avez empêché son retour. À titre d'illustration~: «~[Symptôme] est apparu [à quelle fréquence]. J'ai écarté [deux hypothèses] par des tests. La cause était [cause]. J'ai ajouté [un test ou une alarme].~»

\textbf{Barème.} 2 pour les quatre éléments. 1 pour un récit sans méthode.

\ivround{2}
\ivsays{Concevez un carnet d'ordres pour un instrument. Il y a environ 1\,000 niveaux de prix actifs, et un million de mises à jour par seconde. Comment trouvez-vous un niveau de prix, et comment annulez-vous un ordre~? Un défaut de cache coûte 100 nanosecondes. Pourquoi cela compte-t-il~?}
\ivstop
\textbf{Réponse modèle.} Utilisez un tableau indexé par le prix en ticks, centré sur le prix courant, de sorte qu'une recherche ne prenne qu'une étape. Chaque niveau possède une file d'ordres sous forme de liste chaînée, et une table de hachage qui associe le numéro d'ordre à son nœud rend l'annulation à temps constant. Un arbre équilibré compte environ $\log_2 1\,000 = 10$ étapes. Si chacune est un défaut de cache, cela coûte 1 microseconde, soit tout le budget à un million de mises à jour par seconde. Mutualisez les nœuds, afin de ne pas allouer de mémoire sur ce chemin. Dites ce que vous feriez si les prix sortaient du tableau.

\textbf{Barème.} 2 pour le tableau, l'annulation et l'argument sur le cache. 1 pour un arbre de complexité correcte et sans argument sur le cache.

\ivround{3}
\ivsays{Vous attribuez à 100\,000 ordres un hachage sur 32 bits comme identifiant. Combien de paires en collision attendez-vous~? Quelle est la probabilité d'en avoir au moins une~? Que feriez-vous~?}
\ivstop
\textbf{Réponse modèle~: environ 1.2 paire, et une probabilité d'environ 69~\%.} Le nombre de paires est $n(n - 1)/2 = 5 \times 10^9$, et chacune a une probabilité de $2^{-32}$, soit environ 1.16. La probabilité qu'il n'y en ait aucune est d'environ $e^{-1.16}$, donc il y en a au moins une avec $0.69$. Avec 64 bits, le nombre attendu est d'environ $3 \times 10^{-10}$. Mieux encore, il vaut mieux ne pas hacher~: utiliser un numéro de séquence.

\textbf{Barème.} 2 pour les chiffres et le remède. 1 pour «~le problème des anniversaires~» avec un seul chiffre.

\ivround{4}
\ivsays{Des messages arrivent au hasard à raison de 800\,000 par seconde, et un serveur en traite un million par seconde, avec des temps de service aléatoires. Quel est le temps moyen passé dans le système, et combien de messages s'y trouvent~? Et à 950\,000~?}
\ivstop
\textbf{Réponse modèle~: 5 microsecondes et 4 messages, puis 20 microsecondes et 19.} Pour ce modèle, le temps vaut $1/(\mu - \lambda)$, et le nombre vaut $\lambda$ fois le temps, ce qui est la loi de Little\index{loi de Little}. Passer de 80~\% à 95~\% d'occupation multiplie le délai par quatre, et les rafales aggravent les choses. Il faut donc dimensionner pour les rafales, et surveiller le taux d'occupation comme une alarme précoce.

\textbf{Barème.} 2 pour les deux, avec la raison de la hausse brutale. 1 pour un seul.

\ivround{5}
\ivsays{Un million de messages de 64 octets par seconde pendant 6.5 heures par jour. Combien par jour~? Pour une année de 250 jours, compressée au rapport 5 pour 1~? À quelle vitesse les disques doivent-ils être pour rejouer à dix fois la vitesse~?}
\ivstop
\textbf{Réponse modèle~: 1.5 téraoctet, environ 75 téraoctets, et 128 mégaoctets par seconde depuis le disque, ou 640 s'il est stocké brut.} $64 \times 10^6 \times 23\,400 = 1.5 \times 10^{12}$. Diviser par 5 puis multiplier par 250 donne 75 téraoctets. La relecture demande $64 \times 10 = 640$ mégaoctets par seconde après décompression, donc 128 depuis le disque au rapport 5 pour 1, et la décompression elle-même doit suivre à 640, ce qui est souvent le goulot d'étranglement.

\textbf{Barème.} 2 pour les trois. 1 pour deux.

\ivround{6}
\ivsays{Deux threads incrémentent chacun un compteur partagé un million de fois, sans verrou. Où peut-il finir~? Comment le corrigeriez-vous~? Je donne ensuite à chaque thread son propre compteur, dans la même structure, et c'est lent. Pourquoi~?}
\ivstop
\textbf{Réponse modèle~: n'importe où entre 2 et 2 millions, et faux partage.} Un incrément est une lecture, une addition et une écriture. Les entrelacements peuvent perdre presque tous ces incréments. Il est surprenant que 2 soit possible~: un thread peut conserver une ancienne valeur pendant presque tout le travail de l'autre, et cela deux fois. Le correctif consiste à utiliser un incrément atomique, ou des compteurs par thread additionnés plus tard. Deux compteurs sur la même ligne de cache de 64 octets obligent les cœurs à se renvoyer la ligne~: il faut les placer sur des lignes différentes.

\textbf{Barème.} 2 pour la course, une réponse inférieure à 2 millions avec une raison, et le faux partage. 1 pour deux de ces éléments.

\ivround{7}
\ivsays{Le graphique des latences montre des valeurs typiques de 20 microsecondes, mais un trader affirme qu'à chaque minute le système se bloque pendant environ 50 millisecondes. Votre 99.9\up{e}~centile paraît correct. Comment est-ce possible~? Comment le trouvez-vous~?}
\ivstop
\textbf{Réponse modèle.} Un blocage de 50 millisecondes par minute représente $0.05/60 = 0.083~\%$ du temps, donc il dépasse le 99.9\up{e}~centile en temps. En nombre d'échantillons, c'est beaucoup moins si le test envoie peu de messages pendant ce temps. Tracer la latence dans le temps, avec le maximum par seconde. L'aligner sur les tâches planifiées~: rotation des journaux, minuteur, compaction de la mémoire, instantané. Une période régulière est un indice fort. Modifier une chose à la fois, et vérifier que le blocage a disparu.

\textbf{Barème.} 2 pour le calcul des centiles et la méthode. 1 pour une liste de causes sans test.

\ivround{8}
\ivsays{Que voulez-vous me demander~?}
\ivstop
\textbf{Réponse modèle.} Demander comment le code arrive en production, ce qui se passe après une panne, et ce qu'un ingénieur qui arrive livre pendant son premier mois.

\textbf{Barème.} 2 pour de vraies questions sur le fonctionnement du travail. 1 pour des questions générales.

\medskip
\textbf{Lecture du total.} Ces seuils sont des repères, et non une note de passage fixée par une entreprise. Un total de 13 à 16 serait solide. Pour 9 à 12, il faut retravailler les tours perdus. En dessous de 9, il faut d'abord reprendre les chapitres correspondants.

\iqwrong{Répondre uniquement par la complexité d'une structure de données. Ici, un facteur constant fait en général toute la différence, et il faut dire comment on le mesurerait.}
\iqwhy{Le logiciel de trading exige à la fois la justesse et la vitesse. Un vrai entretien comporterait aussi du codage en direct, que les livres de codage de cette série entraînent.}
\end{solution}

\begin{solution}{iq:iv2:mock-interviews:8}
Il faut garder masqué tout ce qui se trouve sous cette ligne, et faire glisser la feuille vers le bas d'une ligne d'arrêt à la fois. Si la réponse diffère du modèle, il faut poursuivre à partir du modèle.

\ivround{1}
\ivsays{Quelle est la chose la plus rapide que vous ayez jamais construite, et comment savez-vous à quelle vitesse elle allait~?}
\ivstop
\textbf{Une bonne réponse} donne un chiffre, accompagné de la manière dont il a été mesuré. À titre d'illustration~: «~[Système] a traité [débit] avec un 99\up{e}~centile de [valeur], mesuré [comment, de quel point à quel point]. L'étape la plus lente était [étape]. J'examinerais ensuite [quoi].~»

\textbf{Barème.} 2 pour une queue mesurée et la manière de la mesurer. 1 pour une moyenne.

\ivround{2}
\ivsays{Votre processeur tourne à 3 gigahertz, et vous disposez d'une microseconde entre le paquet et l'ordre. Combien de cycles cela représente-t-il~? Un défaut de cache vers la mémoire coûte 100 nanosecondes. Combien de défauts pouvez-vous vous permettre~?}
\ivstop
\textbf{Réponse modèle~: 3\,000 cycles, et au plus 10 défauts, sans rien d'autre.} Un défaut vaut 300 cycles. En pratique, on ne peut en accepter que très peu sur le chemin critique, si bien que l'objectif de conception est d'avoir les données prêtes dans le cache, et de toucher peu de lignes.

\textbf{Barème.} 2 pour les chiffres et la conclusion de conception. 1 pour les chiffres seuls.

\ivround{3}
\ivsays{Une trame occupe 84 octets sur le fil, en comptant tout. Combien de temps faut-il pour l'émettre à 10 gigabits par seconde~? La lumière se propage dans la fibre à deux tiers de sa vitesse dans l'air. Deux bourses sont distantes de 1\,200 kilomètres. Quel est le délai à sens unique par fibre, puis par l'air~? Qu'est-ce que cela indique quant à l'endroit où investir vos efforts~?}
\ivstop
\textbf{Réponse modèle~: 67 nanosecondes, 6 millisecondes et 4 millisecondes.} $84 \times 8/10^{10} = 67.2$ nanosecondes. À 200\,000 kilomètres par seconde, 1\,200 kilomètres prennent 6 millisecondes, et à 300\,000, 4. Sur cette liaison, un écart de 2 millisecondes écrase donc tout travail logiciel. À l'intérieur d'un même site, en revanche, le fil est court et votre propre code compte. L'effort va là où se joue la course.

\textbf{Barème.} 2 pour les trois valeurs et la conclusion. 1 pour deux valeurs.

\ivround{4}
\ivsays{Un message traverse cinq étapes indépendantes. Chacune dépasse 10 microsecondes dans 1~\% des cas. Quelle part des messages subit au moins une étape lente~? Et si un message attend la plus lente de 20 passerelles en parallèle~? Et si les moments lents coïncident~?}
\ivstop
\textbf{Réponse modèle~: 4.9~\%, 18.2~\%, et elle tend vers 1~\%.} $1 - 0.99^5 = 4.9~\%$, et $1 - 0.99^{20} = 18.2~\%$. La queue d'une chaîne est donc bien pire que celle d'une étape, et la diffusion en parallèle la multiplie. Si les moments lents surviennent tous ensemble, par exemple à cause d'un blocage commun, la part des messages lents tend vers 1~\%, mais toutes les étapes sont alors lentes en même temps. Mesurer de bout en bout, et non étape par étape.

\textbf{Barème.} 2 pour les deux valeurs et le cas corrélé. 1 pour une seule valeur.

\ivround{5}
\ivsays{Une branche est mal prédite dans 20~\% des cas, à 15 cycles chacune. Quel est son coût moyen~? Vous parcourez un champ de 8 octets sur 1\,000 structures de 64 octets. Combien de lignes de cache touchez-vous~? Combien si le champ se trouve dans son propre tableau~?}
\ivstop
\textbf{Réponse modèle~: 3 cycles, soit 1 nanoseconde~; 1\,000 lignes contre 125.} $0.2 \times 15 = 3$. Un tableau de champs regroupe 8 valeurs par ligne de 64 octets, soit 125 lignes. Si les données sont froides, cela fait 100 microsecondes contre 12.5. Ici, la disposition en mémoire vaut plus que la branche. Un bon candidat ajoute qu'une branche supprimée peut aussi donner un temps d'exécution plus stable, ce qui intéresse la queue.

\textbf{Barème.} 2 pour les deux. 1 pour une seule.

\ivround{6}
\ivsays{Pourquoi contourner le noyau~? Supposez qu'un appel système coûte environ 1 microseconde dans chaque sens. Que perdez-vous~?}
\ivstop
\textbf{Réponse modèle.} Un envoi et une réception retirent environ 2 microsecondes du chemin, soit le double du budget du tour 2. Vous évitez en outre les interruptions et l'ordonnancement, qui pèsent plus encore sur la queue. En contrepartie, on fait tourner un cœur à pleine puissance, et l'on perd les protections et les outils du noyau~: on assume les bogues de la pile réseau, sa supervision et sa reprise sur incident. Cela en vaut la peine là où la course se joue à la microseconde, et c'est excessif là où ce n'est pas le cas.

\textbf{Barème.} 2 pour le chiffre, la queue et un coût réel. 1 pour la seule vitesse.

\ivround{7}
\ivsays{Vous avez apporté une modification qui, selon vous, fait gagner 2 nanosecondes sur un message qui en prend 500, avec un écart type de 50. Combien d'échantillons par version vous faut-il à un seuil de 5~\% et une puissance de 80~\%~? Et si les échantillons successifs ont une autocorrélation de 0.9~? Comment conduisez-vous cette mesure, et que peut-il vous tromper~?}
\ivstop
\textbf{Réponse modèle~: environ 9\,800, puis environ 186\,000.} $2 \times (1.96 + 0.84)^2 \times 50^2/2^2 = 9\,800$. Si les échantillons suivent un AR(1), l'autocorrélation multiplie la variance d'une moyenne par $(1 + \rho)/(1 - \rho) = 19$. Il faut alterner les mesures, sur le même cœur et dans le même passage, car des passages entiers diffèrent davantage que cet effet. Ce qui peut tromper~: la fréquence du processeur, une disposition mémoire qui change avec le code, la montée en température, et une moyenne qui bouge alors que la queue ne bouge pas. Comparer les centiles qui importent, et figer la configuration. Si c'est la disposition du code qui a changé, il faut alors essayer plusieurs dispositions.

\textbf{Barème.} 2 pour le chiffre, la corrélation et l'alternance. 1 pour le chiffre seul.

\ivround{8}
\ivsays{Que voulez-vous me demander~?}
\ivstop
\textbf{Réponse modèle.} Demander quel est le budget actuel et où il part, comment on saurait que la modification a aidé le P\&L, et comment le matériel et le logiciel travaillent ensemble ici.

\textbf{Barème.} 2 pour de vraies questions sur la mesure du succès. 1 pour des questions générales.

\medskip
\textbf{Lecture de votre total.} Ces fourchettes sont des conseils, et non un seuil de réussite fixé par une firme. Un total de 13 à 16 serait solide. Pour 9 à 12, il faut retravailler les tours perdus. En dessous de 9, il faut d'abord reprendre les chapitres correspondants.

\iqwrong{Annoncer une moyenne faible, ou un indice de référence que vous ne savez pas expliquer. La queue et la façon dont vous l'avez mesurée sont ce que ce recruteur veut entendre.}
\iqwhy{Le travail sur la latence n'a d'intérêt que si sa mesure est fiable. Un vrai recruteur demanderait quel matériel vous avez réellement utilisé. Ses chiffres différeraient de ces valeurs rondes, il faut donc le demander.}
\end{solution}

\begin{solution}{iq:iv2:mock-interviews:9}
Garder couvert tout ce qui se trouve sous cette ligne, et faire descendre la feuille d'une ligne d'arrêt à la fois. Si votre réponse diffère du modèle, reprendre à partir de celui-ci.

\ivround{1}
\ivsays{Pourquoi le risque, plutôt que le trading~?}
\ivstop
\textbf{Une réponse solide} est positive à l'égard du travail. À titre d'illustration~: «~J'aime me demander ce qui peut mal tourner sur l'ensemble d'un livre, et pouvoir dire non avec des chiffres. Dans [exemple], j'ai [trouvé ou mesuré quelque chose], et [ce qui a changé].~»

\textbf{Barème.} 2 pour une raison positive avec un exemple. 1 pour une raison seule. 0 si cela sonne comme un second choix.

\ivround{2}
\ivsays{Un livre a un écart type quotidien de 10 millions de dollars. Quelle est sa VaR à un jour à 99~\%~? Et à 95~\%~? Et sa VaR à dix jours à 99~\%~? Quelle hypothèse venez-vous d'utiliser, et dans quels cas échoue-t-elle~?}
\ivstop
\textbf{Réponse modèle~: 23.3 millions, 16.4 millions et 73.6 millions de dollars.} Les multiplicateurs valent 2.326 et 1.645. Dix jours donnent $\sqrt{10} = 3.16$, en supposant des jours indépendants et identiquement distribués, et un livre inchangé. Cela échoue avec le regroupement de volatilité, avec les tendances, et quand on ne peut pas sortir en dix jours.

\textbf{Barème.} 2 pour les chiffres et l'hypothèse. 1 pour les chiffres seuls.

\ivround{3}
\ivsays{Deux desks, avec des écarts types de 10 et 20 millions de dollars, et une corrélation de 0.3. Quel est le total~? Quelle part du risque apporte chacun~? En crise, la corrélation passe à 1. Quelle est alors la VaR totale~?}
\ivstop
\textbf{Réponse modèle~: environ 24.9 millions, 26~\% et 74~\%, et environ 69.8 millions de dollars de VaR.} La variance est $100 + 400 + 2 \times 0.3 \times 10 \times 20 = 620$. La part de chaque desk dans la variance est sa propre variance plus sa moitié du terme croisé~: $160/620$ et $460/620$. Avec une corrélation de 1, le total vaut 30, et $30 \times 2.326 = 69.8$, contre 57.9 aujourd'hui. La diversification ne vaut donc 12 millions de dollars de VaR que tant qu'elle dure.

\textbf{Barème.} 2 pour les chiffres et le point de crise. 1 pour le total seul.

\ivround{4}
\ivsays{Sur 250 jours, à 99~\%, combien d'exceptions attendriez-vous~? Vous en avez constaté 8. Quelle en est la vraisemblance si le modèle est juste~? De combien la VaR est-elle trop basse~? Que pourriez-vous vérifier d'autre~?}
\ivstop
\textbf{Réponse modèle~: 2.5, environ 0.4~\%, et environ 26~\%.} Le décompte suit une loi binomiale de moyenne 2.5, et la probabilité d'en observer 8 ou plus est d'environ 0.4~\%. La queue des données vaut $8/250 = 3.2~\%$, et $z = 1.852$. Donc si la forme est normale, l'écart type devrait être $2.326/1.852 = 1.26$ fois plus grand. Il faut vérifier si les exceptions se regroupent, si le P\&L de la position était hypothétique ou réel, et si elles se sont concentrées sur un seul desk.

\textbf{Barème.} 2 pour la moyenne, une probabilité de la bonne taille et le facteur d'échelle. 1 pour la moyenne et «~trop nombreuses~».

\ivround{5}
\ivsays{Comparez l'expected shortfall à 97.5~\% avec la VaR à 99~\%. Et si les rendements sont normaux~? Et s'ils suivent une loi de Student $t$ à 4 degrés de liberté, mise à l'échelle pour le même écart type~?}
\ivstop
\textbf{Réponse modèle~: 2.34 contre 2.33 écarts types, et environ 2.82 contre 2.65.} Pour une loi normale, $\varphi(1.96)/0.025 = 2.34$, presque identique. Pour la loi de Student $t$, elles divergent~: la VaR à 99~\% vaut 2.65, et l'expected shortfall 2.82. Un expected shortfall mesure à quel point les pertes sont graves au-delà du seuil de coupure, donc il réagit aux queues d'une façon qu'un quantile ne fait pas. Il est aussi plus difficile à soumettre à un backtest.

\textbf{Barème.} 2 pour les chiffres de la loi normale et la différence conceptuelle. 1 pour la seule différence.

\ivround{6}
\ivsays{Une position vaut 100 millions de dollars. L'action s'échange pour 20 millions par jour, et vous vous autorisez 20~\% de ce volume. Combien de temps faut-il pour sortir~? Si vous vendez régulièrement, de quel facteur votre risque pendant la sortie dépasse-t-il le risque d'une journée sur la position entière~?}
\ivstop
\textbf{Réponse modèle~: 25 jours, et environ 3 fois.} $100/(0.2 \times 20) = 25$. En vendant régulièrement, vous détenez $k/25$ de la position les jours successifs, si bien que la variance vaut $\sum_{k = 1}^{25}(k/25)^2 = 8.84$, et la racine carrée vaut 2.97. Une VaR sur un horizon d'un jour sous-estime ce risque d'un facteur voisin de 3. En situation de stress, le volume baisse et tout le monde vend, donc la sortie est encore plus longue.

\textbf{Barème.} 2 pour les jours et le facteur. 1 pour les jours.

\ivround{7}
\ivsays{Une exposition de 100 millions de dollars, avec une probabilité de défaut de 2~\% dans l'année et une perte de 60~\% en cas de défaut. Quelle est la perte attendue~? Quel spread de crédit cela correspond-il, approximativement~?}
\ivstop
\textbf{Réponse modèle~: 1.2 million de dollars, et environ 120 points de base.} $0.02 \times 0.6 \times 100$. Le spread est à peu près égal à la probabilité de défaut multipliée par la perte. Les spreads réels sont en général plus élevés, car ils rémunèrent aussi le risque et la liquidité.

\textbf{Barème.} 2 pour les deux, et une raison pour un écart avec le marché. 1 pour un seul.

\ivround{8}
\ivsays{Un trader de premier plan dépasse sa limite de 30~\%. Il affirme que c'est une excellente opération et qu'il repassera sous la limite dès demain. Son gérant est d'accord. Que faites-vous~? Ensuite, que voulez-vous me demander~?}
\ivstop
\textbf{Réponse modèle.} Consigner le dépassement, et le remonter par la procédure habituelle, comme pour n'importe qui. Demander que la transaction soit ramenée dans les limites, ou qu'une limite temporaire soit approuvée par la personne compétente, par écrit. Évaluer ce que la position perdrait en scénario de stress. Un dépassement peut être acceptable~; un dépassement non signalé ne l'est jamais. Rester courtois. Puis demander à quelle fréquence les limites sont modifiées, et pourquoi, et quand le risque a arrêté une opération pour la dernière fois à raison.

\textbf{Barème.} 2 pour avoir suivi la procédure et demandé un chiffre de stress, sans hostilité. 1 pour l'un des deux. 0 pour avoir laissé passer parce que le trader est bon.

\medskip
\textbf{Lecture de votre total.} Ces fourchettes sont des conseils, et non un seuil de réussite fixé par une firme. Un total de 13 à 16 serait solide. Pour 9 à 12, il faut retravailler les tours perdus. En dessous de 9, il faut d'abord reprendre les chapitres correspondants.

\iqwrong{Réciter un cadre par cœur. Un responsable des risques veut entendre ce qu'une mesure suppose, quand elle cesse de fonctionner et ce que vous feriez le jour venu.}
\iqwhy{Les responsables des risques ont besoin de chiffres auxquels on croit, et ils doivent savoir tenir bon. Un vrai recruteur demanderait aussi quelle réglementation s'applique à sa propre banque, et elle évolue~; il faut donc consulter le texte en vigueur.}
\end{solution}

\begin{solution}{iq:iv2:mock-interviews:10}
Garder couvert tout ce qui se trouve sous cette ligne, et faire descendre la feuille d'une ligne d'arrêt à la fois. Si votre réponse diffère du modèle, reprendre à partir de celui-ci.

\ivround{1}
\ivsays{Dites-moi ce que vous avez gagné, ce que vous avez perdu, et ce que vous feriez différemment. Sans nommer d'anciens employeurs, et sans rien de confidentiel. Je veux savoir comment vous raisonnez.}
\ivstop
\textbf{Une réponse solide} est honnête sur la pire perte, et sur la part des résultats qui tient à la compétence. À titre d'illustration, avec vos propres éléments~: «~Sur [période], j'ai géré [type de risque] à [volatilité]. Mon pire mois a été [chiffre], à cause de [raison]. J'ai ensuite modifié [quoi]. Je peux vous montrer [ce qui est autorisé], et je ne peux pas partager [ce qui ne l'est pas].~»

\textbf{Barème.} 2 pour une vraie perte, un changement et une ligne claire sur la confidentialité. 1 pour des chiffres sans perte. 0 pour avoir dévoilé les secrets de quelqu'un d'autre.

\ivround{2}
\ivsays{Vous achetez 100 de spot, et vous vendez à découvert 100 d'un perpetual swap avec 25 de marge. Le financement que vous recevez est de 15~\% par an, et il coûte 0.1~\% par côté, soit quatre côtés au total, pour ouvrir et clore. Quelle est votre rentabilité sur capital pour un an~? Le cash rapporte 5~\%. Quel financement fait que cette opération ne rapporte pas davantage que le cash~?}
\ivstop
\textbf{Réponse modèle~: environ 11.7~\%, et environ 6.65~\%.} Le capital vaut 125. Le revenu net est $15 - 0.4 = 14.6$, ce qui, rapporté aux 125, donne 11.68~\%. En posant $(f - 0.4)/125 = 5~\%$, on trouve $f = 6.65$. C'est donc une opération de revenu, reposant sur un financement qui peut chuter vite, et les frais comptent pour une détention courte.

\textbf{Barème.} 2 pour la rentabilité sur l'ensemble du capital et le seuil de rentabilité. 1 pour la seule mention des 15~\%.

\ivround{3}
\ivsays{Votre perpetual vendu à découvert est liquidé si le cours monte de 25~\%, puisque vous disposez de 25 de marge. La volatilité est de 80~\% par an, sans dérive. Quelle est la probabilité de toucher ce niveau en 30 jours~? Et avec 50 de marge~? Et quelle est la rentabilité avec 50~?}
\ivstop
\textbf{Réponse modèle~: environ un tiers, environ 8~\%, et 10~\% sur le même financement avant frais.} La distance logarithmique est $\ln 1.25 = 0.223$, et $\sigma\sqrt{T} = 0.8\sqrt{30/365} = 0.229$. Par le principe de réflexion, la probabilité vaut $2\Phi(-0.97) = 0.33$. Avec 50 de marge, $\ln 1.5 = 0.405$, et $2\Phi(-1.77) = 0.08$. La rentabilité tombe à $15/150 = 10~\%$. Le risque ici est la liquidation d'une opération couverte~: vous perdriez la position vendue tout en détenant le spot, et vous vous retrouveriez alors long au plus haut. Il faut prévoir de recharger la marge à partir de votre spot, et le faire vite, sur une plateforme en tension.

\textbf{Barème.} 2 pour l'argument de réflexion et le compromis. 1 pour «~environ un tiers~» fondé sur l'intuition. Un candidat solide précise que le suivi en continu le surestime légèrement.

\ivround{4}
\ivsays{Chaque plateforme a 2~\% de chances de défaillir sur un an, et vous perdriez tout ce que vous y détenez. Vous détenez 125 dans une seule plateforme. Quelle est votre perte attendue~? Et si vous répartissez entre trois plateformes indépendantes~? Et si vous gardez le spot hors des plateformes et ne répartissez que la marge~?}
\ivstop
\textbf{Réponse modèle~: 2.5, toujours 2.5 avec un pire cas plus petit, et 0.5.} $0.02 \times 125 = 2.5$. Répartir ne change pas la moyenne~: trois plateformes à $41.7$ chacune ont le même total, mais la pire perte probable tombe à un tiers, et vous perdriez quelque chose dans $1 - 0.98^3 = 5.9~\%$ des cas. Garder le spot en conservation propre ramène à $0.02 \times 25 = 0.5$. La décision la plus puissante porte donc sur ce qui est placé sur une plateforme, tout simplement. Notez que ces défaillances ont tendance à être corrélées. Le 2~\% est une estimation, il faut donc le dire.

\textbf{Barème.} 2 pour avoir vu que la répartition réduit la queue et non la moyenne, et pour avoir déplacé le spot. 1 pour l'un des deux.

\ivround{5}
\ivsays{La plateforme A affiche un cours acheteur de 100.20, et la plateforme B un cours vendeur de 100.00. Les frais sont de 5 points de base par côté. Rééquilibrer l'inventaire coûte 8 points de base, et vous ne le payez que si les deux jambes s'exécutent. Votre seconde jambe s'exécute dans 80~\% des cas. Un échec vous coûte au total 15 points de base et ne déplace aucun inventaire entre plateformes. Cela vaut-il la peine~?}
\ivstop
\textbf{Réponse modèle~: non, avec ces chiffres.} Le brut est de 20 points de base, et 10 après deux frais. Le retour de l'inventaire coûte 8, il reste 2. Avec l'échec~: $0.8 \times (10 - 8) - 0.2 \times 15 = -1.4$, car l'inventaire ne doit revenir que lorsque les deux jambes s'exécutent. Ce qui ressemble à de l'argent gratuit est donc une perte, à moins de réduire le coût de rééquilibrage, de le compenser avec les flux en sens inverse, ou d'améliorer le taux d'exécution. Le taux d'exécution au seuil de rentabilité est $15/17$, soit environ 88~\%.

\textbf{Barème.} 2 pour avoir inclus les échecs et le rééquilibrage, et pour avoir nommé un levier. 1 pour 10 points de base sans rééquilibrage.

\ivround{6}
\ivsays{Votre livre a un ratio de Sharpe de 2, et une nouvelle stratégie un ratio de Sharpe de 1.5, sans corrélation. Vous prenez 20~\% de volatilité totale. Comment répartissez-vous le risque~? Quels sont le ratio de Sharpe combiné et le rendement~? Et si le vrai ratio de Sharpe de la nouvelle stratégie est 0.75~?}
\ivstop
\textbf{Réponse modèle~: 16~\% et 12~\% de volatilité, un ratio de Sharpe de 2.5, et un rendement de 50~\%. À 0.75, le ratio de Sharpe obtenu est d'environ 2.05.} Sans corrélation, le risque optimal est proportionnel au ratio de Sharpe, soit $2:1.5$. Donc $20 \times 2/2.5 = 16$, et $20 \times 1.5/2.5 = 12$, avec $\sqrt{16^2 + 12^2} = 20$. Le rendement vaut $2 \times 16 + 1.5 \times 12 = 50$. Si le vrai ratio de Sharpe est 0.75 pour cette répartition, le rendement devient $32 + 9 = 41$, et $41/20 = 2.05$. Je donnerais donc à la nouvelle stratégie moins que ne l'indique la formule, car un ratio de Sharpe mesuré est en général trop élevé, et sa corrélation augmente en période de stress.

\textbf{Barème.} 2 pour la répartition, et pour avoir vérifié qu'on peut se tromper. 1 pour la répartition seule.

\ivround{7}
\ivsays{Vous êtes à ce niveau de risque, et vous avez perdu 6~\% cette semaine. La stratégie est-elle cassée~? À peu près à quelle fréquence cela arrive-t-il~? Que faites-vous lundi~?}
\ivstop
\textbf{Réponse modèle~: pas sur ce seul constat. C'est environ 2.5 écarts types, soit une fois tous les trois ans environ.} La volatilité hebdomadaire vaut $20/\sqrt{52} = 2.77$, et le gain espéré $50/52 = 0.96$. Donc $z = (-6 - 0.96)/2.77 = -2.51$. La probabilité est d'environ 0.6~\% par semaine, soit une fois tous les 3 ans environ. Les queues de la crypto sont plus épaisses, donc l'événement pourrait être plus fréquent.

Le lundi, répartir la perte par stratégie, par plateforme et par facteur. S'agit-il du financement, des clôtures ou des couvertures qui ont échoué, ou d'une position surdimensionnée~? Vérifier que votre risque est bien ce que vous croyez. Réduire le risque si la mécanique a changé, ou si une limite le prescrit. Ne pas réduire seulement à cause d'une perte, et ne pas renforcer pour revenir à l'équilibre. Prévenir son responsable avant qu'il ne pose la question.

\textbf{Barème.} 2 pour la rareté, assortie d'une réserve, et pour une démarche. 1 pour l'un des deux.

\ivround{8}
\ivsays{Si vous étiez moi, pourquoi ne vous embaucheriez-vous pas~? Et que voulez-vous me demander~?}
\ivstop
\textbf{Réponse modèle.} Citer une vraie faiblesse qui compte ici, et ce qu'on fait pour y remédier. À titre d'illustration~: «~Mon historique porte sur [un régime], et je n'ai pas géré [cette taille]. Je voudrais une montée en charge jusqu'à [des éléments probants].~» Poser ensuite les questions que seul un senior peut trancher~: comment il a dimensionné la dernière plateforme défaillante, quel capital la maison engage en cas de drawdown, et comment fonctionne la rémunération en cas de perte.

\textbf{Barème.} 2 pour une vraie faiblesse et des questions sur le capital et le risque. 1 pour l'un des deux. 0 pour «~je travaille trop~».

\medskip
\textbf{Lecture de votre total.} Ces seuils sont des conseils, et non une note de passage d'une quelconque société. Un total de 13 à 16 serait solide pour un dernier tour. Pour 9 à 12, retravaillez les tours que vous avez perdus. En dessous de 9, reprenez d'abord les chapitres correspondants.

\iqwrong{Ne montrer que le rendement. Sur ce marché, ce qui met fin à la carrière d'un trader, c'est le close-out, la défaillance d'une plateforme, ou une taille que personne ne pouvait sortir, et un responsable du trading est attentif à ces points.}
\iqwhy{Un dernier tour vérifie si vous seriez en sécurité avec leur capital. Un véritable responsable du trading connaîtrait votre marché mieux que ce script, et testerait les parties de votre parcours que vous avez omises.}
\end{solution}
